\section{A common family on the prescribed middle interval}

Write $\mathcal C(S)=\{0,\ldots,q-1\}^S$ for the configurations on a
set of sites $S$. Regional tensor products list the complement first.
The following are supporting geometric and counting results for
\cite[Proposition~5.7]{Arad2016RigorousRG}.

\begin{theorem}[Bounds for the selected and prescribed intervals]
    \label{thm:common_middle_interval_bounds}
    \lean{TNLean.MPS.AreaLaw.selected_middle_fixed_interval_bounds}
    \leanok
    Let $L,a,b,\ell$ be nonnegative integers with $a+2\ell\le b$ and
    $b+2\ell\le L-1$, and let $0\le p,s<2\ell$. Put
    $c=a+p+1$, $d=b+s+1$, $u=a+\ell$, and $v=b+\ell+1$. Then
    $c\le d\le L$, $u\le v\le L$, and $\max(c,u)<\min(d,v)$.
\end{theorem}
\begin{proof}\leanok
    The bounds $p,s<2\ell$, together with the two separation inequalities,
    give each endpoint comparison.
\end{proof}

\begin{theorem}[Low-frequency cardinality on the prescribed interval]
    \label{thm:common_middle_low_frequency_cardinality}
    \lean{TNLean.MPS.AreaLaw.actual_common_fixed_middle_mpo_card_low_frequency_le}
    \leanok
    \uses{thm:selected_common_natural_cardinality,
        thm:selected_common_natural_low_frequency_cardinality}
    Retain the geometric data and positions above. Set $E=[a+p+1,b+s+1)$,
    $Z=[a+2\ell+1,b)$, and $\rho=|\mathcal C(E\setminus Z)|$.
    Let $\mathcal W_{k,i,j}$ be the two-marker index set with $q^2$ labels
    for each marker family. Suppose $q\ge2$, $k\ge1$, and
    $i,j\le\lfloor k/(2\ell)\rfloor$. For every nonnegative integer $n$,
    \begin{align}
        n\le |\mathcal W_{k,i,j}|(k+1)\rho^2q^{4\ell}
        \quad\Longrightarrow\quad
        n\le(qk)^{24\ell+4\lfloor k/\ell\rfloor}.
        \label{eq:common_middle_low_frequency_count}
    \end{align}
\end{theorem}
\begin{proof}\leanok
    The chosen position implies $\ell\ge1$. The selected-family count and its
    low-frequency consequence give
    $|\mathcal W_{k,i,j}|(k+1)\rho^2\le(qk)^{20\ell+4\lfloor k/\ell\rfloor}$.
    Multiplication by $q^{4\ell}\le(qk)^{4\ell}$ proves
    (\ref{eq:common_middle_low_frequency_count}).
\end{proof}

\begin{theorem}[A common prescribed-middle family for finite sampling nodes]
    \label{thm:common_middle_grid_mpo_family}
    \lean{TNLean.MPS.AreaLaw.exists_common_fixed_middle_grid_mpo_family}
    \leanok
    \uses{thm:common_middle_interval_bounds,thm:selected_common_natural_mpo_family,
        thm:regional_middle_from_numbered_three_regions,
        thm:interval_mpo_common_family_transfer,
        thm:interval_trace_common_family_transfer,def:regional_middle_transferred_family}
    Retain the geometric data and positions above, and suppose $q\ge1$.
    Let $B,k,i,j$ be arbitrary nonnegative integers. Supply window operators
    $W^L_r,W^R_r$, supported on $[a+r,a+r+2)$ and
    $[b+r,b+r+2)$,
    respectively. Supply site-dependent core tensors of virtual dimension $B$
    on $[a+2\ell,b+1)$ and an arbitrary closing matrix. Let $P_M$ be their
    trace contraction, canonically placed on the full chain.

    Let $\zeta$ be a finite set of sampling nodes. Use the labels
    $L_r,R_r,B_0,B_1,B_2$ and set
    $O=\{L_r,R_r,B_0,B_1,B_2\}\setminus\{L_p,R_s\}$.
    Supply complex weights $z_{\nu,t}$ for $\nu\in\zeta$ and all labels $t$.
    Put $E=[a+p+1,b+s+1)$,
    $J=[a+\ell,b+\ell+1)$,
    $Z=[a+2\ell+1,b)$, and $\rho=|\mathcal C(E\setminus Z)|$. Define
    $Q=|\mathcal W_{k,i,j}|(k+1)\rho^2$.

    There exist nonnegative integers $n_\nu$, degrees $0\le d_{\nu,r}\le k$, and matrices
    $V_{\nu,r}$ on the natural coordinates of $E$, with $0\le r<n_\nu$.
    Read these matrices on $\mathcal C(E)$ using the increasing enumeration
    of $E$. Let $A=E\setminus J$ and $B'=J\setminus E$, and define
    \begin{align}
        I_\nu                   & =\{0,\ldots,n_\nu-1\}\times\mathcal C(A)^2
        \times\mathcal C(B')^2, &
        I                       & =\coprod_{\nu\in\zeta} I_\nu.
        \label{eq:common_middle_node_indices}
    \end{align}
    For $\sigma,\tau\in\mathcal C(J)$, prescribe the middle matrices by
    \begin{align}
        W_{\nu,r,\alpha_r,\alpha_c,\eta_r,\eta_c}(\sigma,\tau)
         & =\mathbf1_{\sigma|_{B'}=\eta_r}\mathbf1_{\tau|_{B'}=\eta_c}
        V_{\nu,r}(\alpha_r\sqcup\sigma|_{E\cap J},
        \alpha_c\sqcup\tau|_{E\cap J}).
        \label{eq:common_middle_node_matrices}
    \end{align}
    The disjoint unions in (\ref{eq:common_middle_node_matrices}) combine
    configurations on $A$ and $E\cap J$ into configurations on $E$.
    The resulting family satisfies $|I|\le|\zeta|Qq^{4\ell}$.
    Each $W_\omega$ has an open matrix product representation on $J$, on the
    paired physical alphabet in increasing site order, with internal bond
    dimensions at most $(\max(1,B))^{2k}$.
    In addition, for every $\omega\in I$ there are a nonnegative integer
    $\chi_\omega\le(\max(1,B))^k$, a closing matrix
    $C_\omega\in\C^{\chi_\omega\times\chi_\omega}$, and site-dependent matrices
    $A_{\omega,v}(x,y)\in\C^{\chi_\omega\times\chi_\omega}$ such that
    \begin{align}
        W_\omega(\sigma,\tau)
         & =\operatorname{tr}\left(C_\omega
        \prod_{v=0}^{|J|-1}A_{\omega,v}(\sigma_{j_v},\tau_{j_v})\right),
        \label{eq:common_middle_trace_presentation}
    \end{align}
    where $j_0<\cdots<j_{|J|-1}$ enumerate $J$ and
    $\sigma,\tau\in\mathcal C(J)$.
    The dimension, closing matrix and tensors are chosen for the same
    $W_\omega$ before the exterior operators and scalar coefficients.

    These choices precede every supported exterior pair $P_L,P_R$ on
    $[0,a+1)$ and
    $[b+2\ell,L)$, respectively. For such a pair, set
    $F_{L_r}=W^L_r$, $F_{R_r}=W^R_r$,
    $F_{B_0}=P_L$, $F_{B_1}=P_M$, $F_{B_2}=P_R$, and
    $D_\nu=\sum_{t\in O}z_{\nu,t}F_t$.
    For every scalar array $c_{\nu,h}$, with $\nu\in\zeta$ and $0\le h\le k$,
    there are joint exterior matrices $M_\omega$ on $\mathcal C(J^c)$ such that
    \begin{align}
        \sum_{\nu\in\zeta}\sum_{h=0}^{k}
        c_{\nu,h}[x^iy^j](D_\nu+xW^L_p+yW^R_s)^h
         & =\left(\sum_{\omega\in I}M_\omega\otimes W_\omega\right)^{\mathrm{phys}}.
        \label{eq:common_middle_node_decomposition}
    \end{align}
    The superscript uses the canonical restriction of a full configuration
    to $J^c$ and $J$; the formal variables commute with matrix coefficients.
    Empty sampling sets, zero virtual dimension, and degree zero are allowed.
\end{theorem}
\begin{proof}\leanok
    \uses{thm:selected_common_natural_mpo_family,
        thm:regional_middle_from_numbered_three_regions,
        thm:interval_mpo_common_family_transfer,
        thm:interval_trace_common_family_transfer,thm:common_middle_interval_bounds,
        thm:chain_mps_of_cut_rank,thm:chain_cut_rank_of_mps}
    Choose the selected natural family separately at each supplied node.
    Pass its joint exterior decomposition to regional coordinates, then
    transfer the middle from $E$ to $J$. The geometric inequalities give the
    required overlap. Moreover $|E\setminus J|+|J\setminus E|\le2\ell$,
    so the canonical family (\ref{eq:common_middle_node_matrices}) has at most
    $Qq^{4\ell}$ members at each node. Taking the disjoint union
    (\ref{eq:common_middle_node_indices}) gives the asserted total count.

    A member of degree $t\le k$ first has open bond bound
    $\max(1,(B^t)^2)$. Since
    $\max(1,(B^t)^2)\le(\max(1,B))^{2k}$, the internal-cut-rank criterion
    gives the uniform bound. The same member is a trace contraction of
    stacked tensors of dimension $B^t$. Restricting a trace contraction to
    fixed configurations outside a subinterval replaces its closing matrix
    $C$ by $SCP$, where $P$ and $S$ are the ordered products before and after
    the subinterval. Inserting fixed row and column configurations at the
    additional sites uses the letters
    $\mathbf1_{x=\eta_v}\mathbf1_{y=\theta_v}\Id_{B^t}$.
    Thus the prescribed transfer leaves the virtual dimension $B^t$ unchanged,
    and $B^t\le(\max(1,B))^k$ gives
    (\ref{eq:common_middle_trace_presentation}).
    These operations apply to the matrices already prescribed in
    (\ref{eq:common_middle_node_matrices}).

    Finally choose the exterior matrices at each
    node for the scalar family $c_{\nu,h}$ and sum their decompositions.
    The finite sum over the disjoint union is exactly
    (\ref{eq:common_middle_node_decomposition}).
\end{proof}
